% page 279 du fac-simile (image p-278 du scan)
% bloc 165.1 x 246.3 mm ; marge gauche 36.9 mm
\pgc{15.240mm}{0.000mm}{
\l{\cel{54}271}
\l{}
\l{\sou{kataklismo}. Subverso di la surfaco di la Terglobo, per inundo,}
\l{\cel{3}sismo (tertremo), e c. - DEFIS.}
\l{}
\l{\sou{katakombo}. Chambrego subsula qua uzesis kom sepulteyo, osteyo ;}
\l{\cel{3}exkavajo adube onu asemblis osti. - DEFIRS.}
\l{}
\l{\sou{katakrezo}. Figuro retorikala, per qua la vorto qua indikas pro-}
\l{\cel{3}pre ta o ca objekto, uzesas por indikar altra objekto qua}
\l{\cel{3}esas sur parto analoga ye la unesma. - DEFIS.}
\l{}
\l{\sou{katalazo}. (\sou{biol.}) Fermento solvebla, deskovrita en la tabako-}
\l{\cel{3}folii, e qua efikas reduktive.}
\l{}
\l{\sou{katalepsio}. (\sou{patol.}) Supreseso semblanta di vivo per la suspens-}
\l{\cel{3}eso di la sentiveso extera e di la moviveso volata kun}
\l{\cel{3}rigidesko kadavrala. - DEFIRS.}
\l{}
\l{\sou{katalizar}.\cel{2}(\sou{trans}.) (\sou{kemio}). Efikigar ta korpi qui, per lia}
\l{\cel{3}nura prezenteso (e sen esar modifikata) eventigas en altra}
\l{\cel{3}korpi modifikesi. - DEF.}
\l{}
\l{katalogo.\cel{2}Listo, indexatra, ek la peci di kolektajo (libri,}
\l{\cel{3}pikturi, graburi,medalii, e c.) - DEFIRS.}
\l{}
\l{\sou{katalpo}. (\sou{bot.)}\cel{2}Arboro ek la familio \textquotedbl{}bignonacei\textquotedbl{}, di qua la}
\l{\cel{3}folii esas larja ed ovala, la flori blanka e reda-puntizita,}
\l{\cel{3}e quan onu kultivas enEuropa kom plezuro-planto. - L.\cel{1}\sou{catalpa}.}
\l{\cel{3}- DEFIS.}
\l{}
\l{kataplasmo.\cel{3}Topiko ek substanco moligiva, forme di pado kom-}
\l{\cel{3}pakta. - DEFIRS.}
\l{}
\l{katapulto.\cel{3}(\sou{che la antiqui)}. Milito-mashino, en qua trabo, per}
\l{\cel{3}efikar quale risorto, lansis projektili tre pezoza. - DEFIRS.}
\l{}
\l{kataro.\cel{2}(\sou{patol.}) Inflameso di mukozo, kun sekreco. - DEFIRS.}
\l{}
\l{katarakto. I. (geol.) Sorto de aquo-fali qui interruptas la}
\l{\cel{3}fluo di fluvio. - II. (\sou{oftalm.}) Opakeso di la kristalino, qua}
\l{\cel{3}interceptas la lumo-radii.-\cel{2}DEFIRS.}
\l{}
\l{katastrofo. I. Inversesko bruska di fortuno. - II. Eventajo}
\l{\cel{3}decidigiva qua solvas la problemo di poemo dramata, epika.}
\l{\cel{3}- DEFIRS.}
\l{}
\l{katedro.\cel{2}Sidilo lokizita alte, e de qua (en kirko, templo,}
\l{\cel{3}doceyo) onu su turnas docoze a sua askoltanti. - DFIRS.}
\l{}
\l{katedralo.\cel{2}Granda kirko di la arkitekturo kristana mezepokala.}
\l{\cel{3}- DEFIRS.}
\l{}
\l{kategorio. I. (filoz.) Singla ek la atributi generala di la}
\l{\cel{4}ento, segun Aristoteles. - II. Singla ek la koncepti aprioria}
\l{\cel{4}di la kompreno homala, segun qui lu konceptas necese la ob-}
\l{\cel{4}jekti di la experienco. - DEFIRS.}
}
